\section*{第 \ref{ch:g10:chemical-species} 章——化学物种：天然与合成}

\begin{solution}{exo:g10:chemical-species:1}
天然：香兰素取自豆荚，咖啡因取自咖啡豆。合成：
阿司匹林，工厂合成的香兰素。
\end{solution}

\begin{solution}{exo:g10:chemical-species:2}
二者属于同一种\omterm{def:g6:pure-substances-and-mixtures:species}{化学物种}，
同一\omterm{def:g6:pure-substances-and-mixtures:species}{化学物种}无论来源如何，性质都相同。
\end{solution}

\begin{solution}{exo:g10:chemical-species:3}
$R_f = 2.4/6.0 = 0.40$.
\end{solution}

\begin{solution}{exo:g10:chemical-species:4}
它使蒸气冷却并重新变为液体。从下端进入的水
充满整个夹套（并最后接触温度最高的蒸气，
冷却效果最好）。
\end{solution}

\begin{solution}{exo:g10:chemical-species:5}
应能比水更好地\omterm{def:g2:mixing-and-dissolving:dissolve}{溶解}该\omterm{def:g6:pure-substances-and-mixtures:species}{化学物种}；与水不互溶；
危害尽量小，且易于除去。
\end{solution}

\begin{solution}{exo:g10:chemical-species:6}
乙醇与水互溶：不会分层，无法分离。
环己烷与水不互溶，且能很好地\omterm{def:g2:mixing-and-dissolving:dissolve}{溶解}芳樟醇。
\end{solution}

\begin{solution}{exo:g10:chemical-species:7}
水在上层：\qty{1}{mL} 二氯甲烷的质量为 \qty{1.33}{g}，
比 \qty{1}{mL} 水的质量更大，因此下沉。
\end{solution}

\begin{solution}{exo:g10:chemical-species:8}
熔化温度偏低（阿司匹林：\qty{135}{\celsius}），且
在一个温度范围内熔化：样品不纯，或者不是阿司匹林。
\end{solution}

\begin{solution}{exo:g10:chemical-species:9}
芳樟醇和乙酸芳樟酯（其两个斑点与两个参照物
位于同一高度）。乙酸芳樟酯上升得更高。
\end{solution}

\begin{solution}{exo:g10:chemical-species:10}
\omterm{def:g10:chemical-species:distillation}{水蒸气蒸馏}；向馏出液中加入环己烷并振荡；
分离各层；\omterm{def:g3:separating-mixtures:evaporate}{蒸发}除去环己烷。
\end{solution}

\begin{solution}{exo:g10:chemical-species:11}
铅笔不溶于\omterm{def:g10:chemical-species:tlc}{洗脱剂}；点样线位于\omterm{def:g10:chemical-species:tlc}{洗脱剂}液面之上，
斑点随\omterm{def:g10:chemical-species:tlc}{洗脱剂}沿薄层板上升，而不会溶入展开槽中。
\end{solution}

\begin{solution}{exo:g10:chemical-species:12}
环己烷沸点为 \qty{81}{\celsius}，芳樟醇沸点为 \qty{198}{\celsius}：
缓慢加热可蒸出环己烷，而芳樟醇留在容器中。
\end{solution}

\begin{solution}{exo:g10:chemical-species:13}
样品很可能就是该\omterm{def:g6:pure-substances-and-mixtures:species}{化学物种}（与同一块板上\omterm{def:g10:chemical-species:natural}{天然物种}的 $R_f$ 相同），
并且很可能是纯的（只有一个斑点，熔化温度范围窄）。
\end{solution}

\begin{solution}{exo:g10:chemical-species:14}
$0.40 \times 5.5 = \qty{2.2}{cm}$。另一个斑点位于 $0.75 \times 5.5 =
\qty{4.1}{cm}$（精确到毫米），因此约在 A 上方 \qty{1.9}{cm} 处。
\end{solution}

\begin{solution}{exo:g10:chemical-species:15}
\qty{200}{mL} 水可\omterm{def:g2:mixing-and-dissolving:dissolve}{溶解} $1.6 \times 0.200 = \qty{0.32}{g}$：
加入 \qty{0.20}{g} 时尚未饱和，还能再\omterm{def:g2:mixing-and-dissolving:dissolve}{溶解} \qty{0.12}{g}。
若\omterm{def:g2:mixing-and-dissolving:dissolve}{溶解}的 \qty{0.20}{g} 随水弃去便会损失；
环己烷对芳樟醇的\omterm{def:g2:mixing-and-dissolving:dissolve}{溶解}能力强得多，可将其中大部分重新萃取回来。
\end{solution}

\begin{solution}{pb:g10:chemical-species:1}
\textbf{1.} 水蒸气将花中带香味的\omterm{def:g6:pure-substances-and-mixtures:species}{化学物种}带出，
送入冷凝器。

\textbf{2.} 冷凝器使蒸气冷却为液体；
冷水从下端进入。

\textbf{3.} 不互溶；密度小于水（会漂浮）。

\textbf{4.} 该层很薄，部分精油\omterm{def:g2:mixing-and-dissolving:dissolve}{溶解}于
水中或分散其中，需要通过萃取将其回收。

\textbf{5.} 它能很好地\omterm{def:g2:mixing-and-dissolving:dissolve}{溶解}芳樟醇，与水不互溶，
且沸点较低（\qty{81}{\celsius}），因此易于除去。

\textbf{6.} 不能：乙醇与水互溶，不会分层。

\textbf{7.} 位于上层（每毫升 \qty{0.78}{g}，比水轻）。若改用
二氯甲烷（每毫升 \qty{1.33}{g}），
\omterm{def:g6:solutions-and-solubility:solute}{溶剂}层将位于下层。

\textbf{8.} 附近不得有明火（易燃）；应在通风橱内操作，
并戴手套，收集废液而不将其倒入水池，
（对水生生物有毒）。

\textbf{9.} 环己烷沸点为 \qty{81}{\celsius}，芳樟醇沸点为
\qty{198}{\celsius}：缓慢加热只能除去\omterm{def:g6:solutions-and-solubility:solute}{溶剂}。

\textbf{10.} $R_f(\text{L}) = 2.4/6.0 = 0.40$；$R_f(\text{A}) = 4.5/6.0
= 0.75$。

\textbf{11.} 芳樟醇和乙酸芳樟酯。

\textbf{12.} 不是：它至少含有两种\omterm{def:g6:pure-substances-and-mixtures:species}{化学物种}，是\omterm{def:g6:pure-substances-and-mixtures:mixture}{混合物}。

\textbf{13.} 为使所有斑点在相同条件下展开，
只有这样才能比较它们的高度。

\textbf{14.} 在相同条件下，其 $R_f$ 与芳樟醇的相等：
它很可能就是芳樟醇，与\omterm{def:g10:chemical-species:natural}{天然物种}相同。

\textbf{15.} $1.0 \times 0.90 = \qty{0.90}{g}$。

\textbf{16.} $0.90/100 = \qty{0.90}{\percent}$。

\textbf{17.} $100/0.009 \approx 11\,000$\,g，约需
\qty{11}{kg} 花朵。

\textbf{18.} 它从植物中\omterm{def:g10:chemical-species:extraction}{提取}得到，未发生\omterm{def:g5:irreversible-changes:chemical-change}{化学变化}。但其中的芳樟醇
与合成芳樟醇属于同一\omterm{def:g6:pure-substances-and-mixtures:species}{化学物种}：化学式相同，
性质也相同。

\textbf{19.} 从 \qty{100}{g} 花中得到 \qty{0.90}{g} 精油。
\end{solution}
